\section{Effective integration and continuous laws}\label{sec:integration}

\subsection{Regular reals and simple-law names}

\begin{definition}[Regular rational interval name]\label{def:regular}
A regular real name is a pair of primitive-recursive rational functions
\(\ell(p),u(p)\) such that
\[
 \ell(p)\le u(p),\qquad
 [\ell(p+1),u(p+1)]\subseteq[\ell(p),u(p)],\qquad
 u(p)-\ell(p)\le2^{-p}.
\]
Arithmetic is performed by rational interval operations with extra guard bits.
\end{definition}

\begin{definition}[Effective simple random variable]\label{def:simple-name}
An effective \(L^1\) simple-law name is a primitive-recursive sequence of
finite rational laws \(X^{(p)}\), together with couplings of any two levels and
a modulus
\[
 \E\abs{X^{(p)}-X^{(q)}}\le2^{-p}
 \qquad(q\ge p).
\]
The couplings are finite common refinements: a finite weighted law carrying
both approximants and having the prescribed marginals.
\end{definition}

This presentation is deliberately operational.  It does not quantify over a
set of measurable functions.  Each level is a finite object, and the stated
\(L^1\) inequality is a rational certificate.

\begin{theorem}[Expectation of an \(L^1\) simple-law name]
\label{thm:l1-expect}
The rational sequence \(e_p=\E[X^{(p)}]\) is a regular Cauchy name.  More
precisely, for \(q\ge p\),
\[
 \abs{e_q-e_p}\le2^{-p}.
\]
Thus \(\E[X]\) is constructively defined.  Equivalent names at zero effective
\(L^1\) distance have the same expectation.
\end{theorem}

\begin{proof}
On the supplied common refinement,
\[
 \abs{e_q-e_p}
 =\abs{\E[X^{(q)}-X^{(p)}]}
 \le\E\abs{X^{(q)}-X^{(p)}}
 \le2^{-p}.
\]
The first inequality is the exact finite absolute-expectation inequality.
Intervals centered at \(e_{p+2}\) with radius \(2^{-(p+2)}\), followed by a
nested-hull normalization, give \Cref{def:regular}.  If two names have
arbitrarily small effective \(L^1\) distance, apply the same estimate across a
common refinement.
\end{proof}

\begin{theorem}[Lifted algebra]\label{thm:l1-algebra}
Finite linearity, monotonicity, almost-sure/zero-\(L^1\) congruence, the absolute
expectation bound, and certified Jensen lift from finite laws to effective
\(L^1\) names.  The output modulus is primitive-recursively computed from the
input moduli and coefficient bounds.
\end{theorem}

\begin{proof}
At output precision \(p\), choose a common finite level with enough guard bits.
Apply the corresponding exact finite theorem.  Bound the discrepancy between
that finite result and each limiting name by the \(L^1\) modulus and
\(\abs{\E Z}\le\E\abs Z\).  For a coefficient \(c\), request \(q\) such that
\(\abs c2^{-q}\le2^{-p}\).  Monotonicity is stated with an effective
nonnegative-difference presentation, rather than an uncheckable pointwise
order on completed limits.
\end{proof}

\subsection{Compact rational quadrature}

\begin{definition}[Piecewise Lipschitz integrand]\label{def:lipschitz}
On a rational compact interval \([a,b]\), a piecewise Lipschitz presentation
contains rational breakpoints, a primitive-recursive rational evaluator
\(g(q,p)\) enclosing \(g(q)\) to precision \(2^{-p}\) at rational \(q\), and
rational constants \(L_j\) satisfying
\[
 \abs{g(x)-g(y)}\le L_j\abs{x-y}
\]
on each piece.  Endpoint values are supplied as one-sided names when needed.
\end{definition}

\begin{theorem}[Effective compact integration]\label{thm:compact-integral}
A piecewise Lipschitz presentation has a regular Riemann integral.  On one
piece of length \(T\), a mesh of width \(h\) gives a left-sum error at most
\(LTh\), plus the accumulated evaluation interval error.  Therefore a
primitive-recursive mesh size achieving error \(2^{-p}\) can be chosen.
\end{theorem}

\begin{proof}
For each cell \([x_k,x_{k+1}]\),
\[
 \abs{\int_{x_k}^{x_{k+1}}g(t)\,dt-g(x_k)(x_{k+1}-x_k)}
 \le\int_{x_k}^{x_{k+1}}L\abs{t-x_k}\,dt
 \le Lh(x_{k+1}-x_k).
\]
Summing yields \(LTh\).  A fully arithmetic proof replaces the integral in
this estimate by upper and lower rectangle sums: Lipschitz continuity bounds
the oscillation on a cell by \(Lh\), so the difference of Darboux sums is at
most \(LTh\).  The finite sums are rational and exact.  Split at the finitely
many breakpoints and allocate the error budget among pieces.
\end{proof}

\begin{definition}[Improper integral with an absolute-tail modulus]
\label{def:improper}
A presentation of \(g:\R\to\R\) includes compact piecewise-Lipschitz data and a
primitive-recursive radius \(R(p)\) satisfying
\[
 \int_{\abs x>R(p)}\abs{g(x)}\,dx\le2^{-p}.
\]
Its integral is the regular limit of compact integrals over
\([-R(p),R(p)]\).
\end{definition}

\begin{theorem}[Effective improper integration]\label{thm:improper}
The construction in \Cref{def:improper} produces a regular real independent of
the certified exhaustion.  Addition, scalar multiplication, and effective
changes of variables preserve it with computable moduli.
\end{theorem}

\begin{proof}
Two truncations beyond \(R(p)\) differ by an integral over a subset of the
absolute tail and hence by at most \(2^{-p}\).  Compact quadrature approximates
each truncation.  For addition, split the tail budget; for a change of variable,
transform the compact quadrature and the tail certificate.  Every output
precision therefore reduces to finitely many rational inequalities.
\end{proof}

\subsection{Densities and expectation}

\begin{definition}[Effective density]\label{def:density}
An effective density is a nonnegative integrand \(f\) with an effective
improper integral equal to one.  A function \(g\) is integrable against \(f\)
when \(\abs g f\) has an absolute-tail modulus and a compact quadrature
presentation.  Define
\[
 \E_f[g]=\int_{-\infty}^{\infty}g(x)f(x)\,dx.
\]
The random variable \(X(x)=x\) has mean \(\int xf(x)\,dx\) when
\(\abs x f(x)\) is effectively integrable.
\end{definition}

\begin{theorem}[Agreement with simple-law expectation]\label{thm:density-agreement}
Let \(f\) be an effective density and let \(g\) be effectively integrable
against \(f\).  There is a primitive-recursive sequence of finite
\emph{rational} laws whose values are rational cell representatives for
\(g\), whose successive levels admit certified finite couplings, and whose
\(L^1\) limit has expectation
\[
 \int_{-\infty}^{\infty}g(x)f(x)\,dx.
\]
Thus the density definition and the simple-law definition of expectation agree.
\end{theorem}

\begin{proof}
Fix an output precision \(p\).  Choose a rational radius \(R\) so that both
\(\int_{|x|>R}f\) and \(\int_{|x|>R}|g|f\) are below a small allocated
budget \(\eta\).  On \([-R,R]\), compute a rational bound \(G\) for
\(|g|\), and choose a rational partition so fine that the oscillation of
\(g\) on every cell is at most \(\eta\).  These choices are primitive
recursive from the compact-continuity and tail moduli.

For each cell \(I_j\), compact quadrature gives rational bounds
\(0\le l_j\le\int_{I_j}f\le u_j\) with
\(u_j-l_j\le\eta/J\), where \(J\) is the number of cells.  Put
\(q_j=l_j\) and add one residual atom of rational mass
\(q_*=1-\sum_jq_j\ge0\), represented at value zero.  The resulting masses
are nonnegative rationals summing exactly to one; a common denominator turns
them into natural weights.  Moreover,
\[
 q_*\le \int_{|x|>R}f+\sum_j(u_j-l_j),
\]
so the residual mass is explicitly small.  If \(r_j\in I_j\) is rational,
the expectation of the finite law is \(\sum_jq_jg(r_j)\).  Its difference
from \(\int gf\) is bounded by three finite estimates: cell oscillation,
\(G\sum_j(u_j-l_j)\), and the certified tail \(\int_{|x|>R}|g|f\).
Choosing \(\eta\) with guard bits makes their sum at most \(2^{-p}\).

For coherence between two levels, refine both rational partitions to their
finite common grid.  Match the rational masses greedily inside common cells;
the unmatched mass is bounded by the two quadrature deficits and the two tail
masses.  Transporting matched mass within a cell costs at most the cell
oscillation, while the unmatched compact mass costs at most \(2G\) times its
mass; the tail contribution is controlled by the \(|g|f\) modulus.  This
finite greedy transport is primitive recursive and supplies a rational
coupling with the required \(L^1\) bound.  Hence the levels form an effective
simple-law name.  The expectation estimate above and
\Cref{thm:l1-expect} show that its regular expectation is exactly the improper
integral.
\end{proof}

\begin{theorem}[Effective continuous LOTUS]\label{thm:continuous-lotus}
For an effective density \(f\) and an effectively integrable \(g\),
\[
 \E[g(X)]=\int_{-\infty}^{\infty}g(x)f(x)\,dx.
\]
For a primitive-recursive effective transformation \(Y=h(X)\), the expectation
of \(k(Y)\) is obtained by integrating \(k(h(x))f(x)\), without first
constructing a separate density for \(Y\).
\end{theorem}

\begin{proof}
The first statement is the definition justified by
\Cref{thm:density-agreement}.  For the second, apply the same finite simple-law
approximants with the composed evaluator \(k\circ h\).  This is the continuous
analogue of finite LOTUS and avoids an unnecessary inverse-function theorem.
\end{proof}

\subsection{CDF, Stieltjes, and layer-cake formulas}

\begin{definition}[Effective CDF name]\label{def:cdf}
For a density presentation, define
\[
 F(t)=\int_{-\infty}^{t}f(x)\,dx
\]
with a two-argument modulus controlling the left tail and compact quadrature.
More generally, a CDF/Stieltjes name consists of rational polygonal
approximants with monotonicity, endpoint, and variation certificates.
\end{definition}

\begin{theorem}[Effective integration by parts]\label{thm:stieltjes}
Let \(F\) be a certified CDF/Stieltjes name and let \(g\) be a compactly
supported piecewise \(C^1\) name, or a piecewise linear rational function.
Then
\[
 \int g\,dF=gF\big|_{a}^{b}-\int_a^bF(t)g'(t)\,dt
\]
with a primitive-recursive error modulus.  Passing to improper limits is valid
when the boundary terms and absolute tails have supplied moduli.
\end{theorem}

\begin{proof}
For polygonal approximants, the formula is finite summation by parts:
\[
 \sum_k g_k(F_{k+1}-F_k)
 =g_bF_b-g_aF_a-\sum_kF_{k+1}(g_{k+1}-g_k),
\]
up to a fixed endpoint convention.  This is a finite algebraic identity.
Uniform variation and continuity bounds control replacement of the names by
finer approximants.  The improper statement follows from the tail certificates.
\end{proof}

\begin{theorem}[Nonnegative layer cake]\label{thm:layer-cake}
For an effectively integrable nonnegative random variable,
\[
 \E[X]=\int_0^\infty\Pp(X>t)\,dt.
\]
For an integrable signed variable and any rational \(a\),
\[
 \E[X]-a
 =\int_a^\infty\Pp(X>t)\,dt
  -\int_{-\infty}^{a}\Pp(X\le t)\,dt.
\]
At \(a=\E[X]\), the two areas are equal.
\end{theorem}

\begin{proof}
At a finite rational simple level with denominator \(D\), the first formula is
exactly \Cref{thm:finite-layer-cake}: the integral is a sum of rectangles of
width \(1/D\).  Truncate the variable at a rational height and apply the
finite identity.  The integrability modulus bounds the discarded
\(\E[(X-B)^+]\), which is exactly the corresponding tail area at the simple
level.  Pass to the effective limit.

For signed \(X\), apply the first formula to \((X-a)^+\) and \((a-X)^+\), then
use \(X-a=(X-a)^+-(a-X)^+\).  The events defining the two positive parts are
\(X>t\) above \(a\) and \(X\le t\) below \(a\), with endpoint differences
confined to zero-width integration slices.  Setting \(a=\E[X]\) makes the left
side zero.
\end{proof}

\subsection{Absolute continuity: what is and is not internalized}

Classically, several formulations of absolute continuity are equivalent, and a
Radon--Nikodym theorem produces a density.  A bare existence conclusion is not
computational data.  We therefore use the following witnessed direction.

\begin{proposition}[Consequences of an explicit density witness]
\label{prop:ac-forward}
Given a nonnegative effective density \(f\), the set function
\(\mu([a,b])=\int_a^bf\) is effectively absolutely continuous with respect to
length on rational finite unions of intervals.  Its CDF is an effectively
absolutely continuous regular function, and its almost-everywhere derivative
is represented by the supplied \(f\) at all certified continuity points.
\end{proposition}

\begin{proof}
Given \(\varepsilon=2^{-p}\), first use the improper-tail modulus to choose
\(R\) with \(\int_{|x|>R}f<\varepsilon/2\).  The piecewise-Lipschitz data
on the compact interval \([-R,R]\) compute a rational upper bound \(M>0\)
for \(f\) there: evaluate one rational point on each piece and add its
Lipschitz constant times the piece length, with one guard bit for the evaluator.
If a finite union of intervals has total length below
\(\delta=\varepsilon/(2M)\), then its integral inside \([-R,R]\) is below
\(\varepsilon/2\), while its part outside contributes below the certified
spatial tail \(\varepsilon/2\).  This proves effective absolute continuity
without adding an independent amplitude-tail oracle.  The CDF increment
formula is its defining integral.  At a continuity point with a modulus, the
difference quotient is trapped between local lower and upper values of
\(f\), giving the derivative estimate.
\end{proof}

\begin{remark}[No uniform constructive Radon--Nikodym converse]
From a code that merely asserts an epsilon--delta property, one cannot in
general extract a primitive-recursive density evaluator and all its integration
moduli.  This report therefore does not claim the classical converse without a
witness.  If a density name is supplied as part of the certificate, all usual
forward equivalences become formalizable.
\end{remark}

\subsection{Cauchy truncation ambiguity}

\begin{proposition}[Cauchy expectation is undefined]\label{prop:cauchy-trunc}
For the standard Cauchy density
\(f(x)=1/(\pi(1+x^2))\), the positive and negative first moments both diverge.
The symmetric truncated first moment is zero, whereas truncation over
\([-R,2R]\) tends to \((\log2)/\pi\).  Thus a principal-value convention does
not define the expectation.
\end{proposition}

\begin{proof}
On a finite interval,
\[
 \int \frac{x}{\pi(1+x^2)}\,dx
 =\frac{1}{2\pi}\log(1+x^2).
\]
Symmetric endpoints cancel.  For \([-R,2R]\), the value is
\[
 \frac{1}{2\pi}\log\frac{1+4R^2}{1+R^2}
 \longrightarrow\frac{1}{2\pi}\log4=\frac{\log2}{\pi},
\]
with an explicit rational bound on the logarithmic ratio error.  On the
positive half-line,
\(x/(1+x^2)\ge1/(2x)\) for \(x\ge1\), so the positive moment dominates a
logarithmically divergent integral; symmetry gives divergence of the negative
part.  By \Cref{def:signed-extended}, the expectation is undefined.
\end{proof}
